% Revised Data & Methodology section
% Paste this into new_intro.tex replacing lines 138–260

\section{Data \& Methodology}

Our primary data source is internal individual and business tax records from 2002--2021, provided by the IRS through the Joint Statistical Research Program. We use a 0.2\% weighted sample of individual returns (approximately 300,000 observations per year) and 100\% samples of business and rental property returns. We supplement the tax data with aggregate series from the Financial Accounts, NIPA, and the Survey of Consumer Finances (SCF). Table 1 presents summary statistics.

The starting point is the Distributional National Accounts of PSZ, which we replicate with internal tax records. We add to this our own estimates of individual-level capital gains, estimated via a three-step procedure: (i) link individuals to their specific portfolio holdings, (ii) capitalize income flows to estimate wealth, and (iii) multiply portfolio value by asset-class-specific returns. We extend prior capitalization approaches in four ways: (i) network tracing of partnership ownership to allocate 90\% of partnership assets to ultimate owners, (ii) newly available tax forms for rental real estate, (iii) private business transaction data for heterogeneous valuations, and (iv) heterogeneous dividend yields and portfolio betas from brokerage data.

A key output is the share of capital gains received by group $g$ (e.g., the top 1\%). Denoting group size by $n_g$ and total population by $n$:
\be
\frac{ n_g E[KG_g]}{nE[KG]} = \frac{ n_g E[W_g]E[r_g]}{nE[KG]}
\ee
For our estimates to be unbiased, a sufficient condition is that capitalization factors and return series are unbiased conditional on the grouping variable. While individual-level estimates contain measurement error, group aggregates will be unbiased. We now summarize each asset class; full details are in Appendix~J.

\subsection{Public equities}

Prior work capitalizes dividends with a common yield \citep{piketty2016distributional} or allows heterogeneous yields by fixed-income type \citep{smith2023}. We use account-level brokerage data from Barber \& Odean (2000) to estimate heterogeneous dividend yields \textit{and} portfolio betas by 25 cells of age and income. Dividend yields increase linearly in age, with the 65+ group averaging 20\% higher yields than the under-35 group.

Because the brokerage data spans 1991--1996, we estimate yields \textit{relative} to the market average and apply these ratios to current-year aggregate yields. Individual public equity wealth is then estimated as nonqualified dividends divided by the group-specific yield. Table~NN compares distributional statistics from our heterogeneous yields with the homogeneous-yield approaches of PSZ and SZZ. For capital gain yields, we estimate cell-level portfolio alphas and betas from a CAPM-style regression of portfolio returns on the market, allowing returns to differ by both average excess return and market exposure. We scale aggregate public equity wealth to match the Financial Accounts.

\subsection{Private business wealth}

\textbf{Valuation.} We estimate firm-level enterprise values by capitalizing EBITDA and sales using valuation ratios derived from private business transaction data (Business Valuation Resources/DealStats), the first use of these data for distributional wealth estimation. We estimate capitalization rates by year, two-digit NAICS industry, firm size, and legal form; final enterprise value is the average of sales-based and EBITDA-based valuations. Size is a particularly important predictor: top-decile firms have EV/EBITDA ratios 1.3 higher than median firms. Net financial assets from Schedule~L are added to convert enterprise values to market values. Details on the BVR data and capitalization rate estimation are in Appendix~J and a companion paper \citep{campbell2025value}.

\textbf{Aggregate results.} Our estimates yield \$17T in aggregate private business wealth in 2017, more than \$6T above SZZ and PSZ, who scale to the Financial Accounts. The gap is largest for partnerships (\$8T vs.\ \$4T), because the FA values partnerships at book value and does not apply market multiples---a growing understatement as intangible capital has risen \citep{bhandari2021sweat,corrado2009intangible} and pass-through activity has expanded \citep{smith2019capitalists,smith2022rise}. Concentration is also substantially higher in our estimates, driven by heterogeneous capitalization rates that value larger businesses at higher multiples and by the high concentration of indirectly owned partnership wealth.

\textbf{Ownership allocation.} S-corp and sole proprietorship wealth is allocated directly from Schedule K-1 and Schedule~C returns. Partnership allocation is more complex because partnerships can be owned by other partnerships, trusts, or corporations. We apply a result from the network economics literature. Let $C$ denote the partnership ownership matrix, where $C_{ij}$ is the fraction of partnership $j$ owned by partnership $i$, and let $a$ be the vector of fundamental asset values estimated above. The total market value vector satisfies $V = a + C \cdot V$, yielding:
\be
V = (I - C)^{-1} \cdot a
\label{eq:network}
\ee
We estimate $C$ from the universe of Form 1065-K1 returns. Of \$7T in partnership enterprise value in 2017, \$6T can be traced to a K1 owner, of which \$2.7T flows to individuals; the remainder is held by C-corporations, trusts, foreign owners, and 10\% by unidentified entities. Capital gains for partnerships and S-corporations equal the change in market value minus change in invested capital; for private C-corporations and sole proprietorships, we apply S\&P 500 price indices.

\subsection{Owner-occupied housing}

Following \citet{larrimore2021recent}, we capitalize property tax deductions from Schedule~A using county-level effective tax rates from the ACS and Decennial Census. This procedure captures approximately 80\% of Financial Accounts housing wealth; following PSZ, we allocate the residual to non-itemizers using SCF averages by income decile and age-by-marital-status cell.

The TCJA of 2018 reduced itemization rates from 30\% to 11\%, dropping our coverage from 77\% to 40\% of FA housing wealth. To recover non-itemizers' housing values post-2018, we match current-year addresses to 2017 returns where property taxes were itemized and update values using FHFA house price indices, raising coverage back to 75\%. The remaining 25\% is allocated via SCF averages from post-TCJA survey years.

\subsection{Other asset classes}

Tenant-occupied housing is valued using property tax data from the Lincoln Institute, with returns from Freddie Mac (residential) and CoStar (commercial) price indices. Fixed-income capital gains are excluded, as real fixed-income returns are frequently negative and contribute minimal distributional variation. Pension wealth is capitalized following PSZ, using a weighted combination of taxable pensions, wages, and tax-exempt pension income, with homogeneous Financial Accounts returns.

\begin{table}[H]
\centering
\caption{Summary of Wealth and Capital Gains Estimation Methods}
\label{tab:estimation_methods}
\small
\begin{tabularx}{\textwidth}{lXX}
\toprule
Asset Class & Wealth Capitalization & Capital Gains Return \\
\midrule
Public equity & Het. yields (Barber \& Odean (2000)) & Het. portfolio $\alpha$ and $\beta$ (Barber \& Odean (2000)) \\[6pt]
Private equity & Het. EV/EBITDA and EV/Sales (DealStats) & Het. returns calculated as change in EV minus investment \\[6pt]
Owner-occupied housing & Het. property tax rates (ACS/Census) & Het. returns (FHFA price indices) \\[6pt]
Tenant-occupied housing & Het. property tax rates (Lincoln Institute) & Het. returns (Freddie Mac and CoStar (commercial)) \\[6pt]
Fixed income & Het. yields (SZZ) & Excluded \\[6pt]
Pension & Capitalize 60\% taxable pensions, 30\% wages, 10\% tax-exempt pensions (PSZ) & Homogeneous (Financial Accounts) \\
\bottomrule
\end{tabularx}
\begin{minipage}{0.95\textwidth}
\vspace{0.5em}
\footnotesize
\textit{Notes:} Table summarizes the capitalization method used to estimate wealth and the return method used to estimate capital gains for each asset class. ``Homogeneous return'' denotes that aggregate Financial Accounts capital gains are divided by aggregate wealth to form a common yield applied to all individuals. ``Heterogeneous'' denotes that capitalization factors or returns vary across individuals based on asset or location characteristics. SZZ refers to \citet{smith2023}; PSZ refers to \citet{piketty2016distributional}. Fixed income capital gains are excluded from the analysis. See text for details.
\end{minipage}
\end{table}
